\section{Hurewicz detection and finite realization}
\label{sec:realization}

Fix \(n\geq 1\).  Proposition~\ref{norm:construction} constructed the
classes \(b_0,\ldots,b_n\), the objects \(Y_0,\ldots,Y_{n+1}\), and
the auxiliary algebras \(L_a\).  We now identify the cobordism Hurewicz
image of each \(b_j\).  This identification will show that the successive
cofibers kill the first powers of the coefficients \(x_j\), and hence
the first powers of the Hazewinkel generators after passage to \(\BP\).

Write \(M=(M_p)_p\), where \(M_p=\MU_{(p)}\), and retain the ideals
\(I_j=(x_0,\ldots,x_{j-1})\subseteq (M_p)_*\) of
Section~\ref{sec:cobordism}.  We set \(I_0=0\).  For each prime, let
\(\overline I_j\subseteq\BP_*\) denote the ideal generated by the
images of these coefficients.  Lemma~\ref{cob:regularity} identifies
\(\overline I_j=(p,v_1,\ldots,v_{j-1})\) for \(j\geq 1\), with
\(\overline I_0=0\).

\subsection{The finite diagram and its effect on homology}

Let \(u_j:\one\to Y_j\) be the unit and let
\[
 \mu_j:S(w_j)\otimes Y_j\longrightarrow Y_j,
 \qquad c_j:Y_j\longrightarrow Y_{j+1}
\]
be, respectively, multiplication by \(b_j:S(w_j)\to Y_j\) and its
cofiber map.  The construction supplies the homotopies
\begin{equation}
 \mu_j\circ(\id\otimes u_j)\simeq b_j,
 \qquad c_j\circ u_j\simeq u_{j+1},
 \qquad \cofib(\mu_j)\simeq Y_{j+1},
 \label{real:finite-identities}
\end{equation}
where the last equivalence respects the cofiber map.  Only these maps,
homotopies, and cofiber identifications, for \(0\leq j\leq n\),
will be needed at individual primes.

They form finite diagram data in \(\cC\).  By the description of
\(\cC\) as a filtered colimit of product categories, they have
representatives on a common \(\U\)-large set of primes.  To be
explicit, choose representatives of the finitely many objects and maps,
then representatives of the finitely many homotopies in
Equation~\eqref{real:finite-identities}.  An equivalence of cofibers is
represented together with an inverse and the two inverse homotopies.
Shrinking the prime set finitely many times makes all these choices
valid together.  Calibrate each of the finitely many source spheres
\(S(w_j)\) with \(S^{d_j}\) at the primes.  We obtain spectra
\(Y_{j,p}\), maps \(u_{j,p},b_{j,p},\mu_{j,p},c_{j,p}\), and
cofiber sequences
\begin{equation}
 \Sigma^{d_j}Y_{j,p}\xrightarrow{\mu_{j,p}}Y_{j,p}
 \xrightarrow{c_{j,p}}Y_{j+1,p}.
 \label{real:primewise-cofibers}
\end{equation}
No descent of the full commutative algebra structures on the \(Y_j\)
is involved in this assertion.

The following elementary observation explains why this finite diagram
is enough to compute homology in every degree.

\begin{lemma}\label{real:cofiber-quotient}
Let \(E\) be a ring spectrum whose coefficient ring \(R=\pi_*E\)
is graded commutative and concentrated in even degrees. Let \(J\)
be a homogeneous ideal of \(R\), and let \(u:S\to Y\) induce
the quotient identification \(E_*Y\cong R/J\). Suppose that
\(b:S^d\to Y\) and \(\mu:\Sigma^dY\to Y\), with \(d\) even, satisfy
\(\mu\circ\Sigma^du\simeq b\).  If the \(E\)-Hurewicz image
of \(b\) is \(a\) times the image of a homogeneous element
\(x\in R_d\), where \(a\in R_0^\times\) and multiplication
by \(x\) is injective on \(R/J\), then the
composite \(S\xrightarrow{u}Y\to\cofib(\mu)\) induces
\[
 E_*\cofib(\mu)\cong R/(J,x).
\]
This is an identification of the entire graded module.
\end{lemma}

\begin{proof}
The map on \(E\)-homology induced by any spectrum map is
\(R\)-linear.  Its source here is the shifted cyclic module
\(\Sigma^dR/J\), and the given homotopy says that it takes its
generator to \(a x\).  It is therefore multiplication by \(a x\)
on all homogeneous elements.  This map is injective.  The cofiber
long exact sequence consequently gives, for every integer \(k\),
\[
 0\longrightarrow (R/J)_{k-d}
 \xrightarrow{\,a x\,}(R/J)_k
 \longrightarrow E_k\cofib(\mu)\longrightarrow 0.
\]
The quotient map in this sequence is induced by the cofiber map, so
the resulting identification is induced by the stated map from the
sphere.  Since \(a\) is a unit, its cokernel is \(R/(J,x)\).
\end{proof}

\subsection{Detection before Euler inversion}

For \(0\leq a<n\), abbreviate
\[
 s_a=(w_a,0,0),\qquad i_a=(w_a,0,-1)=s_a-\tau.
\]
In the arguments below, a quotient by \(t\) denotes the underlying
cofiber of multiplication by \(t\).  It agrees with the underlying
module of the algebra quotient by \(t\).  Thus the \(M\)-version
of Equation~\eqref{eq:target-layer} reads
\begin{equation}
 (L_a^M/t)_i\simeq
 h_\Omega(M\otimes Y_{a+1})_{F(i)}
 \qquad(i\in I).
 \label{real:target-layer-M}
\end{equation}
These equivalences respect the unit, the maps from the base layer,
and the reduction of \(g_a\) used to define \(b_{a+1}\).

\begin{lemma}\label{real:t-injectivity}
Suppose that, on a \(\U\)-large set of primes,
\[
 (M_p)_*Y_{a+1,p}\cong (M_p)_*/I_{a+1}
\]
via the unit.  Then multiplication by \(t\) is injective on
\(\pi_0L_a^M\) at every index.  In particular,
\[
 \pi_0(L_a^M)_i\longrightarrow
 \pi_0(L_a^M[t^{-1}])_i
\]
is injective for every \(i\in I\).
\end{lemma}

\begin{proof}
The coefficient ring of \(M_p\) and its quotient by \(I_{a+1}\)
are concentrated in even degrees.  Every degree \(d(F(i))\) is
even.  Equation~\eqref{real:target-layer-M} therefore gives
\(\pi_1(L_a^M/t)_i=0\) for every fixed \(i\).  In the exact
sequence of the cofiber
\[
 (L_a^M)_{i-\tau}\xrightarrow{t}(L_a^M)_i
 \longrightarrow (L_a^M/t)_i,
\]
this vanishing proves the asserted injectivity.  Localization is the
componentwise telescope along multiplication by \(t\).  Homotopy
groups commute with that filtered colimit, and every transition on
\(\pi_0\) is injective, which proves the last assertion.  Notice
that this argument uses only the evenness of the layer; it makes no
evenness assertion about the completed algebra \(L_a^M\).
\end{proof}

For a coefficient \(x_j\), let \(\mathcal P(x_j)\) denote its
raw powered class in \(\pi_0(L_a^M)_{s_j}\): represent \(x_j\)
at each prime by a map \(S^{d_j}\to M_p\), take its permuted
\(p\)-fold smash, multiply in \(M_p\), and map the resulting
class from \(B(M)\) to \(L_a^M\).  This definition is made
before inverting \(t\).  For \(u\in D^\times\), define
\(\mathcal P(u)\in\pi_0(L_a^M)_0\) in the same way using
scalar maps.  We do not assert additivity of these raw classes
before localization.

\begin{lemma}\label{real:raw-powers}
Assume the hypothesis of Lemma~\ref{real:t-injectivity}.  Suppose
also that \((M_p)_*Y_{a,p}\cong (M_p)_*/I_a\) via the unit on
a \(\U\)-large set, and that the \(M\)-Hurewicz image of
\(b_a\) equals \(u x_a\) for some \(u\in D^\times\).
Then, before inverting \(t\),
\begin{equation}
 \begin{aligned}
  \mathcal P(x_j)&=0
     &&\text{in }\pi_0(L_a^M)_{s_j}\quad(0\leq j<a),\\
  t g_a&=\mathcal P(u)\mathcal P(x_a)
     &&\text{in }\pi_0(L_a^M)_{s_a}.
 \end{aligned}
 \label{eq:before-inversion}
\end{equation}
Here \(g_a\) also denotes its image under base change to \(M\),
and \(\mathcal P(u)\) is a unit at index zero.
\end{lemma}

\begin{proof}
Apply Proposition~\ref{norm:recovery} and
Equation~\eqref{eq:recovery}.  The classes \(x_j\) for
\(j<a\) vanish in its source
\(h_\Lambda(M\otimes Y_a)\), and their images in the target are
the localizations of \(\mathcal P(x_j)\).  The same comparison
takes the Hurewicz image of \(b_a\) to \(t g_a\).  Since it is
a map of rational graded algebras and that Hurewicz image is
\(u x_a\), its value is the localization of
\(\mathcal P(u)\mathcal P(x_a)\).  Lemma~\ref{real:t-injectivity}
now proves both equalities before inversion.  No linearity over
\(D\) is required of the powered comparison.

Raw powering is multiplicative on the represented scalar maps:
the product of the permuted powers of two maps is the permuted
power of their product, followed by the commutative multiplication
of \(M_p\).  Apply this to primewise representatives of \(u\)
and \(u^{-1}\).  It gives
\(\mathcal P(u)\mathcal P(u^{-1})=1\) already in \(B(M)\),
and hence in \(L_a^M\).
\end{proof}

\subsection{Putting the cobordism relation at fixed indices}

Equation~\eqref{eq:cobordism} is an identity of Borel cohomology
classes at each prime.  Its ordinary degree varies with the prime,
whereas \(L_a^M\) is indexed by the fixed lattice \(I\).  We
describe the comparison, including its unit ambiguities and its
effect on transfers.

For \(i=(\lambda,\delta,m)\), the representation underlying
\(W_i\) at the prime \(p\) is
\[
 d(\lambda)\rho+d(\delta)-2mV,
 \qquad
 d(F(i))=p\,d(\lambda)+d(\delta)-qm.
\]
This is a virtual complex representation: \(d(\lambda)\) and
\(d(\delta)\) are even, and \(2V\) is the underlying real
representation of \(V\otimes_{\mathbb R}\mathbb C\).
Complex Thom equivalences thus identify
\(W_i\otimes\Tr M_p\) with
\(\Tr S^{d(F(i))}\otimes\Tr M_p\), as \(M_p\)-modules with
naive \(G\)-action.  The equivalences for virtual negative
summands are obtained by dualizing.  For each fixed index these
equivalences give additive maps
\begin{equation}
 \theta_i:
 \prod_{\U}M_p^{-d(F(i))}(BG)
 \longrightarrow\pi_0B(M)_i
 \longrightarrow\pi_0(L_a^M)_i.
 \label{real:borel-comparison}
\end{equation}
At index zero choose the unit identification, so \(\theta_0\)
is a ring map.  The other \(\theta_i\) respect its scalar action.

If two such Thom equivalences are multiplied and compared with the
equivalence at the sum of their indices, the discrepancy is an
automorphism of a suspended free rank-one \(M_p\)-module over
\(BG\).  It is therefore multiplication by a unit of
\(M_p^0(BG)\).  Consequently products under the maps
\(\theta_i\) agree up to images of units under \(\theta_0\).
All comparisons we will use involve a fixed finite list of indices.
Their representatives and the tensor identifications can be chosen
on a common large prime set; a simultaneous primewise realization
of the entire Picard grading is unnecessary.

\begin{lemma}\label{real:thom-comparisons}
Fix \(0\leq a<n\), and suppose that
\((M_p)_*Y_{a+1,p}\cong(M_p)_*/I_{a+1}\) via the unit on a
\(\U\)-large set.  The maps in
Equation~\eqref{real:borel-comparison} have the following properties.
\begin{enumerate}
\item The image \(\theta_\tau(\epsilon)\) is a unit at index
zero times \(t\).
\item For \(j\leq a\), the image
\(\theta_{s_j}(\Pi_p(x_j))\) is a unit at index zero times
\(\mathcal P(x_j)\).
\item Put \(\delta_j=(0,w_j,0)\).  For \(j\leq a\), the
class \(\theta_{\delta_j}(x_j)\) lies in the image of
\[
 t:\pi_0(L_a^M)_{\delta_j-\tau}
 \longrightarrow\pi_0(L_a^M)_{\delta_j}.
\]
\item Set \(X_{a+1}=\theta_{i_a}(x_{a+1})\).  Under
Equation~\eqref{real:target-layer-M}, its reduction modulo \(t\)
is a unit at index zero times the image of the ordinary coefficient
\(x_{a+1}\) in \((M_p)_*Y_{a+1,p}\), at weight \(w_{a+1}\).
\item Every sequence of finite sums of additive transfers from
prime-to-\(p\) subgroups, in the degree assigned to any fixed
index, has zero image under \(\theta_i\).  The number of
summands may depend on \(p\).
\end{enumerate}
\end{lemma}

\begin{proof}
For the first assertion, recall that \(\epsilon\) is the Euler
class of the complex representation
\(V\otimes_{\mathbb R}\mathbb C\), whose underlying real
representation is \(2V\).  Its zero-section map, after the Thom
identification, is the map \(\one\to T^2\) defining \(t\).
Changing the Thom equivalence changes it by a degree-zero unit.
The second assertion follows by undoing the Thom conversion in the
definition of \(\Pi_p\): the source \(W_{s_j}\) is exactly
the permuted sphere \(N(S(w_j))\).  Both constructions then use
the same permuted smash of the coefficient map and multiplication
in \(M_p\).

At \(\delta_j\) the source is the untwisted sphere
\(\Tr S(w_j)\).  The scalar compatibility of
Equation~\eqref{eq:layer}, followed by its extension to
Equation~\eqref{eq:target-layer}, identifies the reduction of
\(\theta_{\delta_j}(x_j)\) with the usual coefficient action
on \(M\otimes Y_{a+1}\), up to the allowed unit.  This action
is zero because \(x_j\in I_{a+1}\).  Exactness of the
cofiber sequence for \(t\) proves the third assertion.

For the fourth assertion, first note that
\[
 F(i_a)=w_{a+1},\qquad
 d(F(i_a))=p d_a+q=d_{a+1}.
\]
It remains to check the coefficient, since equality of these degrees
alone does not identify the two maps.  Recall that
Equation~\eqref{eq:layer} is constructed through \(\Phi(M)\)
using the chosen \(J^\vee\)-linear sphere comparison
\[
 J^\vee\otimes W_i\simeq
 J^\vee\otimes\Tr S(F(i)).
\]
Extend this comparison to \(J^\vee\otimes\Tr M\).  Also
extend the Borel Thom comparison defining \(\theta_i\) along
the unit of \(J^\vee\).  The two are trivializations of the
same free twisted rank-one module.  Their ratio is an invertible
endomorphism, hence a unit in \(\pi_0\Phi(M)_0\).
With the first trivialization, the constant coefficient map
\(S^{d(F(i))}\to M_p\), followed by the \(J^\vee\)-unit,
is precisely the coefficient supplied by the right-hand comparison
of Equation~\eqref{eq:layer}.  The second trivialization therefore
gives that coefficient times the stated unit.  Extending from this
base layer to Equation~\eqref{real:target-layer-M} uses the actual
unit and tower maps, so it preserves this comparison.

Finally, an additive transfer from \(H\leq G\) factors through
an induced object \(\operatorname{Ind}_H^G Z\).  If
\(|H|\) is prime to \(p\), then \(Z\) is a retract of
\(\operatorname{Ind}_{\{1\}}^H\Res_{\{1\}}^H Z\): the
diagonal and sum maps have composite \(|H|\), which is
invertible \(p\)-locally.  Inducing to \(G\) makes the
original induction a retract of a free \(G\)-induction.  At
each prime, factor a finite sum of transfers through the direct sum
of these induced objects.  That entire sum is a retract of one free
induction on the direct sum of the underlying spectra.  Although its
size may grow with \(p\), these sums form a single sequence
object, and their free inductions are among the objects killed by
the defining localization to \(\cE\).  Tensoring with a
representation sphere preserves induction by the projection
formula, and composing with the Thom maps preserves the
factorization.  Thus the whole transfer error has zero image
already in \(B(M)\), before Euler inversion or completion.
\end{proof}

\subsection{The Hurewicz induction}

\begin{proposition}\label{real:hurewicz}
There is a \(\U\)-large set of primes on which the unit maps
induce, simultaneously for \(0\leq j\leq n+1\),
\begin{equation}
 (M_p)_*Y_{j,p}\cong (M_p)_*/I_j,
 \qquad
 \BP_*Y_{j,p}\cong\BP_*/\overline I_j.
 \label{real:homology-quotients}
\end{equation}
For every \(0\leq j\leq n\), there is a scalar
\(u^{(j)}\in D^\times\) such that the \(M\)-Hurewicz
image of \(b_j\) equals \(u^{(j)}x_j\) in
\(\pi_0h_\Omega(M\otimes Y_j)_{w_j}\).
\end{proposition}

\begin{proof}
At \(j=0\), the quotient assertions hold because \(Y_0=\one\),
and the Hurewicz assertion holds because
\(b_0=\varpi u_0\) represents the sequence \(p=x_0\) at
weight \(w_0\).  Lemmas~\ref{cob:regularity} and
\ref{real:cofiber-quotient}, applied to
Equation~\eqref{real:primewise-cofibers}, give both quotient
assertions for \(Y_1\).

Suppose now that \(0\leq a<n\), that the Hurewicz assertions
hold through \(a\), and that the quotient assertions hold
through \(a+1\).  Lemmas~\ref{real:t-injectivity} and
\ref{real:raw-powers} apply.  We will use
Proposition~\ref{cob:coefficient-relation} to identify the reduction
of \(g_a\), which by construction is the Hurewicz image of
\(b_{a+1}\).

Map Equation~\eqref{eq:cobordism} to \(L_a^M\) at total
index \(s_a=i_a+\tau\).  Its total cohomological degree is
\(-p d_a\).  The index assignments for its factors and error
coefficients are
\[
\begin{array}{c|c@{\qquad}c|c}
 \text{class}&\text{index}&\text{class}&\text{index}\\ \hline
 C,C'&0&\Pi_p(x_j)&s_j\\
 \epsilon&\tau&x_j\ (j\leq a)&\delta_j\\
 x_{a+1}&i_a&D_*&s_a-2\tau\\
 D_{jk}&s_a-\delta_j-\delta_k&H_j&s_a-s_j.
\end{array}
\]
These are fixed indices in \(I\).  They have exactly the required
ordinary degrees: for example,
\(d(F(i_a))=d_{a+1}\),
\(d(F(s_a-2\tau))=p d_a+2q\), and
\(d(F(s_a-\delta_j-\delta_k))=p d_a-d_j-d_k\).
Thus each coefficient in the primewise relation defines a sequence
in the source of the appropriate \(\theta_i\).  There is no
requirement that these coefficients themselves have formulas
independent of \(p\).

Lemma~\ref{real:thom-comparisons} and
Equation~\eqref{eq:before-inversion} now account for every term.
The earlier powered coefficients \(\mathcal P(x_j)\),
\(j<a\), vanish.  The \(\epsilon^2D_*\) term is divisible
by \(t^2\).  Each \(x_j\), \(j\leq a\), assigned to
\(\delta_j\), is divisible by \(t\), so every
\(D_{jk}x_jx_k\) term is also divisible by \(t^2\).
The complete transfer error vanishes, including when the number of
its primewise summands grows.  Products under the Thom comparisons
introduce only index-zero units.  Since \(C,C'\) and
\(\mathcal P(u^{(a)})\) are units, the two remaining leading
terms give
\begin{equation}
 t\bigl(U_1g_a-U_2X_{a+1}\bigr)=t^2z
 \quad\text{in }\pi_0(L_a^M)_{s_a},
 \label{real:cancel-equation}
\end{equation}
where \(U_1,U_2\in\pi_0(L_a^M)_0^\times\) and
\(z\in\pi_0(L_a^M)_{s_a-2\tau}\).  After the error
coefficients in Equation~\eqref{eq:cobordism} have been collected,
the list of terms just used has size bounded in terms of \(a\).

Multiplication by \(t\) from index \(i_a\) to index \(s_a\)
is injective by Lemma~\ref{real:t-injectivity}.  Canceling it in
Equation~\eqref{real:cancel-equation} gives
\[
 U_1g_a-U_2X_{a+1}=tz.
\]
Reduce modulo \(t\) and use
Equation~\eqref{real:target-layer-M}.  The reduction of \(g_a\)
is the \(M\)-Hurewicz image of \(b_{a+1}\), because the
construction of the new cycle and its chosen attaching-map path
commutes with the unit from the sphere to \(M\).
Lemma~\ref{real:thom-comparisons} identifies the reduction of
\(X_{a+1}\) with a unit times the ordinary coefficient
\(x_{a+1}\).  We have therefore proved that the Hurewicz image
of \(b_{a+1}\) is a unit at index zero times that coefficient,
at weight \(w_{a+1}\).

This unit can be replaced by a scalar in \(D^\times\).
Indeed, at index zero the actual layer has
\begin{equation}
 \pi_0(L_a^M/t)_0
 \cong\prod_{\U}(M_p)_0Y_{a+1,p}
 \cong\prod_{\U}\F_p
 \cong D/\varpi.
 \label{real:scalar-layer}
\end{equation}
The identifications are ring identifications via the unit: the
quotient assertion identifies the underlying degree-zero module,
and the unit map is already a ring map in the ultraproduct.
Here \(I_{a+1}\) contains \(p\), and its other generators have
positive ordinary degree.  A unit in \(D/\varpi\) has a
representative whose residue at \(p\) is nonzero on a
\(\U\)-large set.  Lift those residues to elements of
\(\Z_{(p)}^\times\), and choose arbitrary units at the
remaining primes.  This gives the required lift in \(D^\times\).
Its action agrees with that of the original layer unit.  Denote the
resulting scalar by \(u^{(a+1)}\).

The resulting Hurewicz equality is an equality in the one specified
mapping group
\[
 \pi_0\map_{\cC}(S(w_{a+1}),M\otimes Y_{a+1})
 \cong
 \prod_{\U}(M_p)_{d_{a+1}}Y_{a+1,p}.
\]
It consequently holds on a \(\U\)-large set at the actual
primes, with scalar representatives in \(\Z_{(p)}^\times\).
Naturality under the multiplicative projection \(M_p\to\BP\)
gives the analogous Hurewicz equality in \(\BP\)-homology.
Lemma~\ref{cob:regularity} makes both coefficient classes
nonzerodivisors in the current quotients.  Applying
Lemma~\ref{real:cofiber-quotient} gives the two quotient assertions
for \(Y_{a+2,p}\), in all degrees at once.

This completes the induction.  At each step the Hurewicz equality
requires only one new large-set restriction, together with the
finitely many comparisons used at that step.  The calculation of
the entire homology module then holds at every prime of that set
by coefficient linearity and the cofiber exact sequence.  Since
\(n\) is fixed, intersecting the finitely many resulting sets
proves the simultaneous assertion.  In particular, no intersection
over all homological degrees, weights, or infinitely many stages
is taken.
\end{proof}

\subsection{Finiteness and proof of the main theorem}

\begin{lemma}\label{real:finite-descent}
The \(p\)-localizations of finite CW spectra are closed under the
finite suspensions and cofibers in
Equation~\eqref{real:primewise-cofibers}.
\end{lemma}

\begin{proof}
Suspensions preserve this class.  For finite CW spectra \(F_0,F_1\),
compactness of \(F_0\), and the expression of localization as a
filtered colimit of multiplication maps, give
\[
 [ (F_0)_{(p)},(F_1)_{(p)} ]
 \cong [F_0,F_1]\otimes\Z_{(p)}.
\]
Thus a map between these localizations has the form \(f/s\),
where \(f:F_0\to F_1\) is a stable map and \(p\nmid s\).
Precomposing or postcomposing with multiplication by \(s\), an
equivalence after localization, identifies its cofiber with
\(\cofib(f)_{(p)}\).  The cofiber of a stable map of finite
CW spectra is again a finite CW spectrum: represent the map after
a finite suspension by a cellular map between finite CW models
and take its finite mapping cone, then desuspend.  This proves the
claim and applies inductively starting from the sphere.
\end{proof}

\begin{proof}[Proof of Theorem~\ref{thm:main}]
For \(n=0\), choose any prime \(p\) and take the cofiber of
\(p:S_{(p)}\to S_{(p)}\).  Multiplication by \(p\) is
injective on \(\BP_*\), so the cofiber exact sequence gives
\(\BP_*X\cong\BP_*/(p)\), via the map from the sphere.
This is a finite \(p\)-local spectrum.

For \(n\geq1\), fix the construction above and choose a prime
in the common \(\U\)-large set of
Proposition~\ref{real:hurewicz}.  Set \(X=Y_{n+1,p}\).
Equations~\eqref{real:primewise-cofibers} and
Lemma~\ref{real:finite-descent} show that \(X\) is the
\(p\)-localization of a finite CW spectrum, and hence is finite
\(p\)-local in the sense of the theorem.  By
Equation~\eqref{real:homology-quotients} and
Lemma~\ref{cob:regularity}, its actual unit map from the
\(p\)-local sphere induces a surjection
\[
 \BP_*\longrightarrow\BP_*X
 \quad\text{with kernel }(p,v_1,\ldots,v_n).
\]

For both cases, the polynomial cooperation algebra
\(\BP_*\BP=\BP_*[t_1,t_2,\ldots]\) is flat over the
coefficient ring \cite[Section~4.4]{HazewinkelIII}. Thus
\(\BP\)-homology of a spectrum carries its natural
\(\BP_*\BP\)-coaction, and the map from the sphere is
a comodule map. Since it is surjective on homology, naturality
determines the coaction on every element of the quotient: it is
exactly the coaction induced from the coefficient comodule
\(\BP_*\).  Thus the displayed isomorphism has the canonical
quotient comodule structure, and its generator is the image of
\(1\in\BP_0\), in degree zero.  This proves the assertion
for every \(n\), with the prime allowed to depend on \(n\).
\end{proof}
